% ─────────────────────────────────────────────────────────────────────────────
% Section VIII — Conclusion
% ─────────────────────────────────────────────────────────────────────────────

\section{Conclusion}
\label{sec:conclusion}

We introduced CRSC, the first formal metric quantifying backdoor persistence in
post-trained LLMs as a function of model scale.
Combining normalized Frobenius drift ($\Delta_\text{hidden}$) with Shannon
entropy change ($\Delta H$), CRSC yields a score in $[0,1]$ (threshold
$\tau = 0.62$) backed by two theoretical guarantees: monotonic non-decrease in
$N$ (Proposition~\ref{prop:monotone}) and concentration-of-measure convergence
to a proxy for ASR (Proposition~\ref{prop:proxy}).

Experimental validation across 72 runs demonstrated F1\,$=\,0.891$, AUC-ROC\,$=\,0.924$,
and $r\,=\,0.847$ between CRSC and $\log_{10}(N)$.
Real-weight validation on 9 LoRA-injected LLaMA-2 checkpoints (7B and 13B)
confirmed the monotonicity prediction, with the 13B style-trigger run
(CRSC\,$=\,0.628$) crossing $\tau$ as the sole \textsc{high}-risk case.
CRSC establishes \emph{security scaling laws} as a research direction: formal
analogs to capability scaling laws that quantify adversarial robustness rather
than benign performance.
Code is available at \url{https://github.com/ufxa/Cyber-Resilience-Scaling-Coefficient}.

\section*{Acknowledgments}

The authors thank the LICA (Laboratório de Inteligência e Cibersegurança Aplicada)
research group at UFRA for computing resources and the SEC365 team for operational
support. This work was supported in part by the CNPq research infrastructure program.
